\begin{abstract}
With terabyte-scale memory capacity
    and memory-intensive workloads,
memory translation has become a major performance bottleneck.
Many novel hardware schemes are
    developed to speed up memory translation, but
    few are experimented with commodity OSes.
A main reason is that memory management in major OSes, like Linux,
    does not have the extensibility
    to empower emerging hardware schemes.
%    either is hardwired to specific hardware architectures
 %   and lack extensibility,
  %  or is not expressive enough to empower emerging architectures.
    
% to support new schemes that cannot fit in tree-based abstractions.

%    ; supporting a new scheme requires extensive
%    modifications of memory management code,
%    including ``architecture-independent'' portion, creating 
%    obstacles to hardware innovations.

We develop EMT, a pragmatic framework atop Linux to empower
    different hardware schemes of memory translation 
    such as radix tree and hash table.
% termed Extensible Memory Translation (EMT).
EMT provides an architecture-neutral interface that
    1) supports diverse memory translation architectures,
    2) enables hardware-specific optimizations,
    3) accommodates modern hardware and OS complexity,
%    3) support modern memory management features, % like huge pages and page-table isolation,
    and 4) has negligible overhead over hardwired implementations. 
% With EMT, new hardware translation schemes can be effectively
%    supported.
We port Linux's memory management onto EMT and show that 
    EMT enables extensibility without sacrificing performance.
We use EMT to implement OS support for ECPT and FPT, two recent experimental translation 
    schemes for fast translation;
    EMT enables us to understand the OS perspective of these architectures
    and further optimize their designs.
\end{abstract}
